\label{sec:definition}

\begin{definition}[Adams Method]
  Given populations $p_1, \ldots, p_n$ and House size $H$,
  the Adams apportionment finds a common divisor $d > 0$ such that:
  \[
    \sum_{i=1}^{n} \max\!\left(1, \left\lceil \frac{p_i}{d} \right\rceil\right) = H,
  \]
  and assigns each state $s_i = \max(1, \lceil p_i/d \rceil)$ seats.
\end{definition}

\noindent
Note that for states with $p_i > 0$, we have $\lceil p_i/d \rceil \geq 1$
as long as $d \leq p_i$.
Since the minimum state population (Wyoming 2020: 576{,}851) is much
larger than any reasonable divisor, the $\max(1, \cdot)$ is binding only
for the constitutional floor guarantee.

\subsection{Priority Queue Formulation}

Adams as a priority queue uses the first-seat priority
$P_i^{\text{Ad}}(0) = p_i$ (effectively infinite, since every state
gets its first seat) and for subsequent seats:
\[
  P_i^{\text{Ad}}(s) = \frac{p_i}{s},
\]
the same formula as Jefferson.
The difference between Adams and Jefferson is not the priority formula
but the initialisation: Adams assigns 1 seat to each state and then
runs the priority queue for the remaining $H - n$ seats;
Jefferson assigns 0 seats initially and runs for $H$ steps.
In practice, since every state must have at least 1 seat
(constitutional floor), the two methods agree for all states above
the constitutional floor, and the floor is active only for
states with exact quotas below 1.
Adams and Jefferson use the same priority formula $p_i/s$
but Adams's base allocation guarantees the constitutional floor
more naturally.

\subsection{Rounding Threshold}

For exact quota $q_i = p_i/d \in (s, s+1)$:
\begin{itemize}
  \item Adams assigns $s_i = s+1$ seats (ceiling).
  \item The threshold for receiving $s+1$ seats is $q_i > s$,
    i.e., any fractional part triggers rounding up.
\end{itemize}
This is the lowest possible threshold: any state with a non-integer
quota receives the upper quota under Adams.
Compare to HH's threshold $\sqrt{s(s+1)} > s$ and Webster's
threshold $s + 0.5 > \sqrt{s(s+1)}$.

\subsection{Worked Example}

For populations $(100, 50, 10)$ and $H = 5$ seats,
exact quotas at $d = 160/5 = 32$ are $(3.125, 1.5625, 0.3125)$.
Adams assigns ceilings $(4, 2, 1)$ = 7 seats, too many.
Adjusting $d$ upward to $d = 36$: quotas $(2.78, 1.39, 0.28)$,
ceilings $(3, 2, 1)$ = 6 seats, still too many.
At $d = 53$: quotas $(1.89, 0.94, 0.19)$,
ceilings $(2, 1, 1)$ = 4 seats, too few.
At $d = 45$: quotas $(2.22, 1.11, 0.22)$,
ceilings $(3, 2, 1)$ = 6 seats.
At $d = 55$: quotas $(1.82, 0.91, 0.18)$,
ceilings $(2, 1, 1)$ = 4 seats.
The Adams divisor is between 45 and 55; binary search converges
to $d \approx 50$ giving quotas $(2.0, 1.0, 0.2)$,
ceilings $(2, 1, 1)$ = 4 seats.
The divisor must be decreased further; at $d = 49$:
quotas $(2.04, 1.02, 0.20)$, ceilings $(3, 2, 1)$ = 6.
At $d = 51$: quotas $(1.96, 0.98, 0.196)$, ceilings $(2, 1, 1)$ = 4.
There is no divisor that gives exactly 5 seats under Adams for this example,
illustrating that Adams may not have a divisor that exactly achieves
a given $H$ --- in practice, the divisor giving the closest total
to $H$ is used, with tie-breaking by population.
